% Informe típico de ingeniería: lo que se lea al compilarlo viaja dentro del APK.
\documentclass[11pt,a4paper]{article}
\usepackage[T1]{fontenc}
\usepackage{lmodern}
\usepackage[spanish]{babel}
\usepackage{csquotes}
\usepackage[margin=2.5cm]{geometry}
\usepackage{amsmath,amssymb,amsthm}
\usepackage{graphicx,xcolor}
\usepackage{booktabs,multirow,float,caption}
\usepackage{enumitem,fancyhdr,titlesec,parskip}
\usepackage{siunitx}
\usepackage{hyperref}
\pagestyle{fancy}\fancyhf{}\rhead{Informe}\cfoot{\thepage}
\newtheorem{teorema}{Teorema}
\title{Informe de laboratorio}\author{TeXsting}\date{\today}
\begin{document}
\maketitle
\tableofcontents
\section{Introducción}\label{sec:intro}
Un \enquote{circuito RC} con $R = \SI{10}{\kilo\ohm}$ y $C=\SI{100}{\nano\farad}$ tiene
\begin{equation}\label{eq:tau}
  \tau = RC = \SI{1}{\milli\second}, \qquad v(t) = V_0\left(1-e^{-t/\tau}\right).
\end{equation}
Ver la ecuación~\eqref{eq:tau} y la sección~\ref{sec:intro}. \textbf{Negrita}, \textit{cursiva}, \textsc{versalitas}, \texttt{código}; acentos: canción, pingüino, ñandú, ¿qué?, ¡sí!
\begin{teorema} Para todo $n\in\mathbb{N}$, $\sum_{k=1}^{n} k = \frac{n(n+1)}{2}$. \end{teorema}
\begin{proof} Inducción sobre $n$. \end{proof}
\subsection{Resultados}
\begin{table}[H]\centering\caption{Mediciones}\label{tab:m}
\begin{tabular}{@{}lS[table-format=2.2]c@{}}\toprule
Medida & {Valor} & Unidad\\\midrule
\multirow{2}{*}{Tensión} & 4.98 & \si{\volt}\\ & 5.02 & \si{\volt}\\
Corriente & 12.30 & \si{\milli\ampere}\\\bottomrule
\end{tabular}\end{table}
\begin{itemize}[noitemsep]\item {\color{red}Rojo} y {\color{blue!60!black}azul}. \item $\alpha,\beta,\gamma,\int_0^\infty f(x)\,dx$ \end{itemize}
\begin{enumerate}[label=\alph*)]\item Uno \item Dos\end{enumerate}
\begin{align}
  \nabla\cdot\mathbf{E} &= \frac{\rho}{\varepsilon_0} & \mathbf{F} &= q(\mathbf{E}+\mathbf{v}\times\mathbf{B})
\end{align}
\url{https://github.com/P3dr1t0/texsting}
\end{document}
